\begin{helper}
Let \(N,n,m\) be natural numbers with \(n\le N\) and \(m\le N\), and let
\(M\subseteq \operatorname{Fin}(N)\) have cardinality \(m\).  Put
\[
q=\frac{m}{N}
\]
using Lean's real-division convention when \(N=0\).  Then, for every
\(L\in\mathbb R\),
\[
( \binom{N}{n})^{-1}
\sum_{S\in \binom{\operatorname{Fin}(N)}{n}} \exp\!\bigl(L\,|S\cap M|\bigr)
\le (1-q+q e^L)^n .
\]
\end{helper}
\begin{proof}
If \(N=0\), then \(n=m=0\).  There is exactly one \(0\)-element subset of
\(\operatorname{Fin}(0)\), namely the empty set, its intersection with \(M\)
has cardinality \(0\), and \(\binom 00=1\).  The left hand side is therefore
\(\exp(0)=1\).  The right hand side is \((1-0+0)^0=1\), so the desired
inequality holds.

Assume \(N>0\).  Write \(U=\operatorname{Fin}(N)\), identify \(U\) with
\(\{1,\ldots,N\}\), and let \(v_i=1\) if \(i\in M\) and \(v_i=0\) otherwise.
Let
\[
b=(\underbrace{1,\ldots,1}_{n\text{ times}},
  \underbrace{0,\ldots,0}_{N-n\text{ times}}).
\]
For a permutation \(\pi\in S_N\), the sum
\(\sum_i b_i v_{\pi(i)}\) counts the marked elements among the first \(n\)
positions of \(\pi\).  Every \(n\)-element subset \(S\subseteq U\) occurs as
the first \(n\) positions of exactly \(n!(N-n)!\) permutations.  Since
\(N!=\binom Nn n!(N-n)!\), the subset average on the left hand side is
exactly
\[
g(b)=\frac{1}{N!}\sum_{\pi\in S_N}
\exp\!\left(L\sum_{i=1}^N b_i v_{\pi(i)}\right).
\]

Now draw \(X=(X_1,\ldots,X_n)\) independently and uniformly from \(U\), and
let \(c_i\) be the number of times \(i\) appears among the \(X_j\).  Then
each \(c_i\) is a nonnegative integer and \(\sum_i c_i=n\).  If
\(Z=\sum_{j=1}^n v_{X_j}\), then \(Z=\sum_i c_i v_i\).  Applying an
independent uniform permutation \(\pi\) to all labels does not change the law
of the sequence \(X\), so for this fixed count vector \(c\) the symmetrized
conditional moment generating function is
\[
g(c)=\frac{1}{N!}\sum_{\pi\in S_N}
\exp\!\left(L\sum_{i=1}^N c_i v_{\pi(i)}\right),
\]
and the with-replacement moment generating function is the average of
\(g(c(X))\) over all sequences \(X\).

Fix one count vector \(c\).  Sort it in nonincreasing order and call the
sorted vector \(x\).  The vector \(b\) is already nonincreasing.  For
\(1\le r\le N\), the sum of the first \(r\) coordinates of \(b\) is
\(\min(r,n)\).  The corresponding sum for \(x\) is at least this number.
Indeed, if \(r\le n\) and \(\sum_{h\le r}x_h<r\), then the sorted
nonnegative integer condition forces \(x_r=0\); all later coordinates are
also \(0\), so the total sum of \(x\) is \(<r\le n\), contradicting
\(\sum_h x_h=n\).  If \(r\ge n\), then no positive coordinate can occur after
position \(r\), since that would give at least \(r+1>n\) positive integer
coordinates and hence total sum at least \(r+1\).  Thus the first \(r\)
coordinates already contain the whole mass \(n\).  Therefore
\[
S_r(x):=\sum_{h=1}^r(x_h-b_h)\ge0\qquad(1\le r\le N),
\]
and \(S_N(x)=0\).

We now prove the finite majorization-to-convex-hull step.  Suppose first
that the current sorted vector \(x\) is not equal to \(b\).  Let \(i\) be the
first index with \(x_i\ne b_i\).  Since all earlier discrepancies are zero,
\(S_i(x)=x_i-b_i\), and the prefix inequality gives \(x_i>b_i\).  Because
the total discrepancy is zero, some later index has negative discrepancy;
let \(j>i\) be the first index with \(x_j<b_j\).  For every
\(i\le r<j\), no negative discrepancy has yet appeared after the positive
discrepancy at \(i\), hence \(S_r(x)>0\).  Define
\[
\delta=\min\Bigl(x_i-b_i,\ b_j-x_j,\ \min_{i\le r<j}S_r(x)\Bigr).
\]
All entries and slacks are integers and each quantity in this minimum is
positive, so \(\delta\) is a positive integer.  Let \(y\) be obtained from
\(x\) by replacing \(x_i\) by \(x_i-\delta\), replacing \(x_j\) by
\(x_j+\delta\), and leaving every other coordinate fixed.  The inequalities
\(\delta\le x_i-b_i\le x_i\) and \(\delta\le b_j-x_j\) show that \(y\) has
nonnegative integer entries.  Its total sum is still \(n\).  Its prefix
slacks are unchanged for \(r<i\), are \(S_r(x)-\delta\) for
\(i\le r<j\), and are unchanged again for \(r\ge j\); hence all prefix
slacks remain nonnegative.

The vector \(y\) need not be sorted.  Let \(y^\downarrow\) be its
nonincreasing rearrangement.  For every \(r\), the first \(r\) entries of
\(y^\downarrow\) have maximal sum among all \(r\)-coordinate selections from
\(y\), so their sum is at least the sum of the first \(r\) entries of the
particular ordering \(y\).  Since those prefix sums of \(y\) dominate the
prefix sums of \(b\), the sorted vector \(y^\downarrow\) also majorizes
\(b\), and its total sum is still \(n\).

The same transfer also gives the required convex-hull inclusion.  Since
\(i<j\), \(b_i\ge b_j\), and from \(x_i>b_i\) and \(x_j<b_j\) we get
\(x_i>x_j\).  Put \(\theta=\delta/(x_i-x_j)\).  The identity
\[
x_i-x_j=(x_i-b_i)+(b_i-b_j)+(b_j-x_j)
\]
has nonnegative terms on the right, with the first and last positive; because
\(\delta\le x_i-b_i\) and \(\delta\le b_j-x_j\), we have
\(0\le\theta\le1\).  On coordinates \(i,j\),
\[
(x_i-\delta,\ x_j+\delta)
=(1-\theta)(x_i,x_j)+\theta(x_j,x_i),
\]
and all other coordinates are unchanged.  Thus \(y\) is a convex combination
of \(x\) and the transposition of \(x\) swapping \(i\) and \(j\).  If
\(\rho\) is a permutation sorting \(y\), then
\[
y^\downarrow=\rho(y)
=(1-\theta)\rho(x)+\theta\,\rho((ij)x),
\]
so \(y^\downarrow\) is still in the convex hull of permutations of \(x\).

It remains to show that the iteration reaches \(b\).  Use the integer
measure
\[
D(x)=\sum_{h=1}^N |x_h-b_h|.
\]
Before resorting, the two affected terms decrease by \(\delta\) each,
because \(x_i>b_i\), \(x_j<b_j\), and
\(\delta\le x_i-b_i,\ b_j-x_j\).  Hence \(D(y)=D(x)-2\delta<D(x)\).
Resorting cannot increase the distance to the already sorted vector \(b\):
an adjacent inverted pair has values \(a\le d\) but appears in the order
\((a,d)\), while the corresponding two entries of \(b\) satisfy
\(b_r\ge b_s\).  The elementary inequality
\[
|a-b_s|+|d-b_r|\le |a-b_r|+|d-b_s|
\]
for \(a\le d\) and \(b_s\le b_r\) says that swapping the larger value leftward
does not increase the two-coordinate contribution to \(D\).  Repeating such
adjacent swaps until the vector is nonincreasing gives
\(D(y^\downarrow)\le D(y)<D(x)\).  Since \(D\) is a nonnegative integer, the
process terminates.  A terminal sorted vector that still majorizes \(b\) must
equal \(b\), because otherwise the first differing index \(i\), the later
index \(j\), and the positive transfer above would still be available.

Starting from the sorted rearrangement \(x=c^\downarrow\), each step replaces
the current sorted vector by a new sorted vector lying in the convex hull of
permutations of the previous one.  Composing these finite convex
representations, and using that \(c^\downarrow\) is itself a permutation of
\(c\), gives weights \(w_\sigma\ge0\), indexed by \(\sigma\in S_N\), whose
sum is \(1\), such that
\[
b=\sum_{\sigma\in S_N} w_\sigma\,\sigma(c).
\]

The function \(g\) is convex in its vector argument: each summand is the
exponential of a linear function, hence convex, and \(g\) is their finite
average with nonnegative coefficient \(1/N!\).  It is also symmetric:
for every permutation \(\sigma\), reindexing the outer sum by
\(\pi\mapsto \pi\circ\sigma\) gives \(g(\sigma(c))=g(c)\).  Jensen's
inequality applied to the displayed convex combination gives
\[
g(b)\le \sum_{\sigma\in S_N} w_\sigma g(\sigma(c))
       = \sum_{\sigma\in S_N} w_\sigma g(c)
       = g(c).
\]
This inequality holds for every fixed sequence \(X\), so averaging over all
with-replacement sequences gives
\[
\mathbb E e^{LY_{\mathrm{without}}}\le
\mathbb E e^{LZ_{\mathrm{with}}}.
\]

For one with-replacement draw, the probability of landing in \(M\) is
\(|M|/|U|=m/N=q\), and the \(n\) draws are independent.  The one-draw moment
generating function is therefore \(1-q+q e^L\), and multiplication over the
\(n\) independent draws gives
\[
\mathbb E e^{LZ}=(1-q+q e^L)^n.
\]
Combining this identity with the domination just proved is exactly the
claimed inequality.
\end{proof}
